% Generated draft: requires mathematical and source review.
\documentclass[11pt]{article}
\usepackage[margin=25mm]{geometry}
\usepackage{fontspec}
\setmainfont{Latin Modern Roman}
\usepackage{amsmath,amssymb,mathtools,unicode-math}
\setmathfont{Latin Modern Math}
\usepackage{enumitem,needspace,xurl,hyperref}
\setlength{\parindent}{0pt}
\setlength{\parskip}{5pt}
\setlength{\emergencystretch}{3em}
\newcommand{\N}{\mathbb{N}}
\newcommand{\Z}{\mathbb{Z}}
\newcommand{\Q}{\mathbb{Q}}
\newcommand{\R}{\mathbb{R}}
\newcommand{\C}{\mathbb{C}}
\newcommand{\F}{\mathbb{F}}
\newcommand{\A}{\mathbb{A}}
\newcommand{\PP}{\mathbb{P}}
\newcommand{\E}{\mathbb{E}}
\newcommand{\eps}{\varepsilon}
\newcommand{\dd}{\,\mathrm{d}}
\begin{document}
\section*{UnsolvedMath ID: 800399}
\textbf{Source problem code:} catalogue entry
\subsection{Definitions and mathematical statement}
Let \(n\geq 1\) be an integer. Consider the families \(f_c(x)=x^2+c\), \(f_c(x)=x^m+c\), or another natural one-parameter family of rational maps on \(\PP^1\), such as \(f_b(x)=bx+x^{-1}\). The source does not specify the range of \(m\), the admissible parameter locus, or the class of other families intended. For a polynomial family of degree at least two, the usual dynatomic polynomial is
\[
\Phi_{n,f}(x,c)=\prod_{d\mid n}\bigl(f_c^{\circ d}(x)-x\bigr)^{\mu(n/d)},
\]
where \(f_c^{\circ d}\) denotes the \(d\)-fold iterate and \(\mu\) is the M\"obius function. A point has formal period \(n\) when it lies on this dynatomic locus; its exact period can be a proper divisor of \(n\) at special parameters. For rational families, the corresponding construction on \(\PP^1\), including the treatment of points at infinity and parameter degenerations, must also be specified.

Under the usual convention, \(Y_1^{\mathrm{dyn}}(n)\) parametrizes a map in the family together with a marked point of formal period \(n\), and \(Y_0^{\mathrm{dyn}}(n)\) is the corresponding curve with a marked cycle, obtained by identifying markings related by iteration. The supplied source instead assigns ``points, respectively cycles'' to \(Y_0^{\mathrm{dyn}}(n)\) and \(Y_1^{\mathrm{dyn}}(n)\), while part (c) associates \(\Phi_n\) with the latter curve. This labeling conflict is unresolved; the requests about both curves are retained.

For each intended family and each \(n\geq1\), address the following questions.
\begin{enumerate}
\item[(a)] Describe the rational primes of bad reduction for each of \(Y_0^{\mathrm{dyn}}(n)\) and \(Y_1^{\mathrm{dyn}}(n)\).
\item[(b)] Give a natural geometric or dynamical explanation for the sporadic primes of bad reduction. The source explicitly describes this terminology as imprecise: if \(S_f\) denotes an intended finite set of bad primes associated with the family, these are bad-reduction primes \(p\) satisfying
\[
p\nmid n\qquad\text{and}\qquad p\notin S_f.
\]
Neither \(S_f\) nor an intrinsic rule defining it is supplied.
\item[(c)] Describe the primes \(p\) for which the reductions of either or both curves are geometrically irreducible. For an intended integral model \(\mathcal Y_i(n)\), this means asking whether
\[
\mathcal Y_i(n)\times_{\operatorname{Spec}\Z}
\operatorname{Spec}\overline{\F}_p
\]
is irreducible, whenever that fiber is defined. The source specifically asks, for \(Y_1^{\mathrm{dyn}}(n)\), for the primes such that the reduced polynomial \(\overline{\Phi}_n(x,z)\) does not factor in \(\overline{\F}_p[x,z]\), where \(z\) denotes the family parameter.
\end{enumerate}

The questions require a specification of the integral models and the meaning of bad reduction, including whether one studies an affine dynatomic model, its normalization, or a smooth projective completion. These choices are not supplied and are not resolved here. Furthermore, geometric irreducibility of a hypersurface as a topological space does not by itself exclude repeated factors in its defining polynomial. Thus the polynomial nonfactorization formulation in part (c) requires additional qualifications concerning reducedness and the chosen presentation before it can be identified with the geometric formulation.
\subsection{Short English statement}
For the indicated one-parameter dynamical families and every positive integer \(n\), describe the bad-reduction primes and the primes of geometric irreducibility of the dynamical curves marking formal-period points or cycles. Explain geometrically or dynamically any bad primes outside the divisors of \(n\) and a finite family-dependent exceptional set. The source leaves the curve labeling, integral models, exceptional set, and some equivalences unresolved.
\subsection{Sources}
\begin{itemize}
\item Local supplied UnsolvedMath registry record, permanent ID \(800399\), exact source code \texttt{catalogue entry}, title: \emph{Open Problems in Arithmetic Dynamics --- Question 22}. The record attributes the question to Benedetto, DeMarco, Ingram, Jones et al., with proposed year 2018.
\item Supplied source: Benedetto, DeMarco, Ingram, Jones et al., \emph{Current Trends and Open Problems in Arithmetic Dynamics} (2018), Question 22, \url{https://arxiv.org/abs/1806.04980}. The supplied provenance records access on 2026-10-16 and extraction by \texttt{arxiv\_latex}; these are supplied provenance claims, not an independently verified visit or transcription.
\item The registry supplies the status \texttt{open}. This is a source label, not an independently verified assessment of the current literature. No separate dated literature assessment is supplied.
\end{itemize}
\end{document}
